% Part: computability
% Chapter: recursive-functions
% Section: non-pr-functions

\documentclass[../../../include/open-logic-section]{subfiles}

\begin{document}

\olfileid{cmp}{rec}{npr}
\olsection{ఆదిమ పునరావృత్తం కాని ప్రమేయాలు}

అంతర్బోధతో గణనీయాలని భావించే ప్రమేయాలన్నీ ఆదిమ పునరావృత్త ప్రమేయాలు
కావు. ఏకస్థానిక ఆదిమ పునరావృత్త ప్రమేయాలన్నిటినీ
$f_0$, $f_1$, $f_2$,~\dots అని జాబితా చేసి, ఇన్‌పుట్~$y$పై $f_x$
విలువను కార్యసాధకంగా గణించవచ్చని అంతర్బోధతో స్పష్టంగా ఉండాలి. మరో
మాటలో, కింది విధంగా నిర్వచించిన ప్రమేయం $g(x,y)$ గణనీయం:
\[
g(x,y) = f_x(y)
\]
అయితే కింది ప్రమేయం కూడా గణనీయమే:
\begin{eqnarray*}
h(x) & = & g(x,x) + 1 \\
& = & f_x(x) +1.
\end{eqnarray*}
ప్రతి ఆదిమ పునరావృత్త ప్రమేయం $f_i$కు, $i$ వద్ద $h$, $f_i$ల విలువలు
వేరుగా ఉంటాయి. కాబట్టి $h$ గణనీయం, కానీ ఆదిమ పునరావృత్త ప్రమేయం
కాదు; $g$ గురించి కూడా ఇదే చెప్పవచ్చు. ఇది కాంటర్ వికర్ణీకరణ వాదనకు
``కార్యసాధక'' రూపం.

గణనీయమై, ఆదిమ పునరావృత్తం కాని ప్రమేయాలకు మరింత స్పష్టమైన ఉదాహరణలు
ఇవ్వవచ్చు. ఉదాహరణకు, $g^n(x)$ అనే సంకేతనం మొత్తం $n$ $g$లతో
$g(g(\dots g(x)))$ను సూచించనివ్వండి; ప్రమేయాల క్రమం
$g_0,g_1,\dots$ను ఇలా నిర్వచించండి:
\begin{eqnarray*}
g_0(x) & = & x+1 \\
g_{n + 1}(x) & = & g_n^x(x).
\end{eqnarray*}
ప్రతి ప్రమేయం $g_n$ ఆదిమ పునరావృత్తమని నిర్ధారించవచ్చు. ప్రతి తరువాతి
ప్రమేయం ముందుదానికన్నా ఎంతో వేగంగా పెరుగుతుంది: $g_1(x)$ అనేది
$2x$కు సమానం, $g_2(x)$ అనేది $2^x\cdot x$కు సమానం; $g_3(x)$
సుమారుగా $x$ సంఖ్యలో $2$లు గల ఘాత
గోపురంలా పెరుగుతుంది. అకెర్మాన్--పీటర్ ప్రమేయం సారాంశంగా
$G(x)=g_x(x)$; ఇది ఏ ఆదిమ పునరావృత్త ప్రమేయంకన్నా వేగంగా పెరుగుతుందని
చూపవచ్చు.

ఆదిమ పునరావృత్త ప్రమేయాలను జాబితా చేసే ప్రశ్నకు తిరిగి వద్దాం. ప్రతి
ఆదిమ పునరావృత్త ప్రమేయానికి సంకేతాత్మక సంకేతనాలను కేటాయించామని
గుర్తుంచుకోండి; కాబట్టి సంకేతనాలను జాబితా చేస్తే సరిపోతుంది. ప్రతి
సంకేతనం $F$కు సహజ సంఖ్య $\#(F)$ను పునరావృత్తంగా ఇలా కేటాయించవచ్చు:
\begin{eqnarray*}
\#(0) & = & \langle 0 \rangle \\
\#(S) & = & \langle 1 \rangle \\
\#(\Proj{n}{i}) & = & \langle 2, n, i \rangle \\
\#(\fn{Comp}_{k,l}[H,G_0,\dots,G_{k-1}]) & = & \langle
3,k,l,\#(H),\#(G_0),\dots,\#(G_{k-1}) \rangle \\
\#(\fn{Rec}_l[G,H]) & = & \langle 4, l, \#(G), \#(H) \rangle.
\end{eqnarray*}
ముందటి విభాగంలోని సంకేతాలను ఉపయోగించి సంఖ్యల ప్రతి క్రమాన్నీ ఒక సహజ
సంఖ్యగా చూడవచ్చనే వాస్తవాన్ని ఇక్కడ వాడుతున్నాం. దీని ఫలితం: ప్రతి
సంకేతానికి ఒక సహజ సంఖ్య కేటాయించబడుతుంది. కొన్ని క్రమాలు (అందువల్ల
కొన్ని సంఖ్యలు) సంకేతనాలకు అనుగుణం కావు. $i$తో సంకేతీకరించిన
సంకేతనం గల ఏకస్థానిక ఆదిమ పునరావృత్త ప్రమేయాన్ని $f_i$గా తీసుకోవచ్చు,
$i$ అటువంటి సంకేతనాన్ని సంకేతీకరిస్తే; లేకపోతే స్థిర $0$ ప్రమేయాన్ని
తీసుకోవచ్చు. ఇలా
ఏకస్థానిక ఆదిమ పునరావృత్త ప్రమేయాలను జాబితా చేయడానికి స్పష్టమైన పద్ధతి
లభిస్తుంది.

(నిజానికి, స్థిర శూన్య ప్రమేయం వంటి కొన్ని ప్రమేయాలు జాబితాలో
ఒకటికన్నా ఎక్కువసార్లు కనిపిస్తాయి. ఇది మన సంకేతీకరణ వల్ల మాత్రమే
కాదు; స్థిర శూన్య ప్రమేయానికి ఒకటికన్నా ఎక్కువ సంకేతనాలు ఉండటం కూడా
కారణం. ఈ పునరుక్తులను గణనీయంగా నివారించలేమని తరువాత చూస్తాం. ఉదాహరణకు,
ఇచ్చిన సంకేతనం స్థిర శూన్య ప్రమేయాన్ని సూచిస్తుందో లేదో నిర్ణయించే
గణనీయ ప్రమేయం లేదు.)

ఇప్పుడు $f_x$ పైన వర్ణించిన జాబితాను సూచించేలా $g(x,y)$ను $f_x(y)$తో
తీసుకోవచ్చు. $g(x,y)$ గణనీయమని మనకు ఎలా తెలుసు? అంతర్బోధతో ఇది
స్పష్టం: $g(x,y)$ను గణించడానికి మొదట $x$ను ``విప్పి'', అది ఏకస్థానిక
ప్రమేయానికి సంకేతనమో కాదో చూడాలి. అయితే, ఇన్‌పుట్~$y$పై ఆ ప్రమేయపు
విలువను గణించాలి.

\begin{tagblock}{TMs}
\begin{digress}
(కొంత శ్రమతో!) ఈ పని చేసే ప్రోగ్రామును Java లేదా C++లో రాయవచ్చని
మీకు ఇప్పటికే నమ్మకం కలిగి ఉండవచ్చు. అప్పుడు, అంతర్బోధతో గణనీయమైన
దేనినైనా ట్యూరింగ్ యంత్రంతో గణించవచ్చని చెప్పే చర్చ్--ట్యూరింగ్
సిద్ధాంతప్రతిపాదనను ఆశ్రయించవచ్చు.

అయితే $g(x,y)$ను గణించే ట్యూరింగ్ యంత్రాన్ని స్పష్టంగా వర్ణించడం,
అది గణనీయమని చూపడానికి మరింత ప్రత్యక్ష పద్ధతి. ముఖ్యంగా, ఇది
చర్చ్--ట్యూరింగ్ సిద్ధాంతప్రతిపాదననూ, అంతర్బోధను ఆశ్రయించడాన్నీ
నివారిస్తుంది. ట్యూరింగ్ యంత్రంపై \emph{అనుకరించగల} గణనా నమూనాను
ఆశ్రయించి $g(x,y)$ గణనీయమని చూపడానికి సరిపడా సాధనాలను త్వరలో
నిర్మిస్తాం; ఆ నమూనానే పునరావృత్త ప్రమేయాలు.
\end{digress}
\end{tagblock}

\end{document}
